% ch14_functors_as_structure_maps.tex — Volume II Chapter 2

\chapter{Functors as Structure Maps}

A functor $F : \mathcal{C} \to \mathcal{D}$ carries objects to objects and
morphisms to morphisms, preserving composition and identity. But what precisely
is being preserved? This chapter examines the structural content of functors —
what they must, may, and cannot preserve — and then extends the picture to
\term{profunctors}, which generalize functors by replacing the target category
with a category of relations.

% ─────────────────────────────────────────────────────────────────────────────
\section{What Functors Preserve and Reflect}

A functor $F : \mathcal{C} \to \mathcal{D}$ \emph{preserves} a property $P$ of
morphisms if: whenever $f$ has property $P$ in $\mathcal{C}$, $F(f)$ has property
$P$ in $\mathcal{D}$. It \emph{reflects} $P$ if: whenever $F(f)$ has $P$ in
$\mathcal{D}$, $f$ has $P$ in $\mathcal{C}$.

\begin{center}
\renewcommand{\arraystretch}{1.4}
\begin{tabular}{lll}
\hline
\textbf{Property} & \textbf{All functors preserve?} & \textbf{Notes} \\
\hline
Isomorphisms & Yes & $F(f)F(f^{-1}) = \operatorname{id}$ \\
Monomorphisms & Not in general & Left adjoints need not preserve monos \\
Epimorphisms & Not in general & Right adjoints need not preserve epis \\
Limits & Right adjoints preserve limits & RAPL \\
Colimits & Left adjoints preserve colimits & LAPC \\
\hline
\end{tabular}
\end{center}

The classification matters: knowing whether a functor is a left or right adjoint
immediately tells you what it preserves, before computing anything specific.

\subsection{Full, Faithful, Essentially Surjective}

\begin{definition}
A functor $F : \mathcal{C} \to \mathcal{D}$ is:
\begin{itemize}
  \item \term{faithful} if each $F_{a,b} : \mathcal{C}(a,b) \to \mathcal{D}(Fa,Fb)$
    is injective;
  \item \term{full} if each $F_{a,b}$ is surjective;
  \item \term{fully faithful} if each $F_{a,b}$ is bijective;
  \item \term{essentially surjective} if for every $d \in \mathcal{D}$, there
    exists $c \in \mathcal{C}$ with $F(c) \cong d$.
\end{itemize}
An equivalence of categories is a functor that is fully faithful and essentially
surjective.
\end{definition}

Fully faithful functors reflect isomorphisms:

\begin{theorem}
If $F : \mathcal{C} \to \mathcal{D}$ is fully faithful and $F(f)$ is an
isomorphism, then $f$ is an isomorphism.
\end{theorem}

\begin{prooflog}
\proofstep{Suppose $F(f) : F(a) \to F(b)$ is an isomorphism in $\mathcal{D}$,
           with inverse $g : F(b) \to F(a)$.}
           {We need to find $h : b \to a$ in $\mathcal{C}$ with $h \circ f = \operatorname{id}_a$
            and $f \circ h = \operatorname{id}_b$.}
\proofstep{Since $F$ is full, there exists $h : b \to a$ in $\mathcal{C}$ with
           $F(h) = g$.}
           {Fullness gives a preimage for every morphism in $\mathcal{D}(F(b), F(a))$.}
\proofstep{$F(h \circ f) = F(h) \circ F(f) = g \circ F(f) = \operatorname{id}_{F(a)} = F(\operatorname{id}_a)$.}
           {Functoriality of $F$ and the fact that $g$ is inverse to $F(f)$.}
\proofstep{Since $F$ is faithful, $F(h \circ f) = F(\operatorname{id}_a)$ implies
           $h \circ f = \operatorname{id}_a$.}
           {Faithfulness: equal images under $F$ imply equal morphisms.}
\proofstep{By the same argument, $f \circ h = \operatorname{id}_b$.}
           {Symmetric argument using $F(f \circ h) = F(f) \circ g = \operatorname{id}_{F(b)}$.}
\proofstep{Therefore $f$ is an isomorphism with inverse $h$.}
           {Both composite conditions are verified.}
\end{prooflog}

% ─────────────────────────────────────────────────────────────────────────────
\section{The Image and Subcategories}

A functor $F : \mathcal{C} \to \mathcal{D}$ has an \term{essential image}: the
full subcategory of $\mathcal{D}$ on objects isomorphic to $F(c)$ for some
$c \in \mathcal{C}$. Full faithfulness means $F$ is an equivalence onto its
essential image.

\begin{definition}
A \term{subcategory} $\mathcal{S} \subseteq \mathcal{D}$ is full if for all
$s, s' \in \mathcal{S}$, $\mathcal{S}(s, s') = \mathcal{D}(s, s')$. The inclusion
functor $\iota : \mathcal{S} \hookrightarrow \mathcal{D}$ is then fully faithful.
\end{definition}

Reflective subcategories arise when the inclusion functor has a left adjoint
(the \term{reflector}):

\begin{definition}
A full subcategory $\iota : \mathcal{S} \hookrightarrow \mathcal{D}$ is
\term{reflective} if $\iota$ has a left adjoint $R : \mathcal{D} \to \mathcal{S}$.
\end{definition}

The reflector $R$ sends each object of $\mathcal{D}$ to the ``nearest'' object
in $\mathcal{S}$. Examples: sheafification (sheaves are a reflective subcategory
of presheaves); abelianization (abelian groups as a reflective subcategory of groups).

% ─────────────────────────────────────────────────────────────────────────────
\begin{nameorigin}
\term{Profunctor} (also \term{distributor}, also \term{bimodule}) was
introduced by \emph{Jean Bénabou} (1973) in his work on bicategories.
The name combines \emph{pro-} (Latin ``in favor of, in place of'') with
\emph{functor}: a profunctor $P : \mathcal{C} \nrightarrow \mathcal{D}$
is ``a functor in place of an ordinary functor'' — generalizing functors
$\mathcal{C} \to \mathcal{D}$ by allowing $\mathbf{Set}$-valued data
$P(c, d)$ for each pair $(c, d) \in \mathcal{C} \times \mathcal{D}$.

The alternative \emph{distributor} (also Bénabou's term) emphasizes
that a profunctor distributes morphisms across the two-sided structure.
\emph{Bimodule} emphasizes the analogy with bimodules over rings: a
profunctor is a $\mathcal{D}$-$\mathcal{C}$-bimodule of sets. All three
names denote the same thing.
\end{nameorigin}

\section{Profunctors}

A functor $F : \mathcal{C} \to \mathcal{D}$ gives a hom-set bijection
$\mathcal{D}(F(c), d)$ for every $c, d$. This is a function of two variables —
contravariant in $c$, covariant in $d$. This is the structure of a profunctor.

\begin{definition}
A \term{profunctor} (also called a \term{distributor} or \term{bimodule})
from $\mathcal{C}$ to $\mathcal{D}$, written $P : \mathcal{C} \nrightarrow \mathcal{D}$,
is a functor
\[
  P : \mathcal{D}^{op} \times \mathcal{C} \;\to\; \mathbf{Set}.
\]
That is, a presheaf on $\mathcal{D}^{op} \times \mathcal{C}$.
\end{definition}

\begin{howtosay}
$P : \mathcal{C} \nrightarrow \mathcal{D}$ is read ``$P$ is a profunctor from
$\mathcal{C}$ to $\mathcal{D}$.'' The slashed arrow $\nrightarrow$ distinguishes
profunctors from ordinary functors.
\end{howtosay}

Every functor $F : \mathcal{C} \to \mathcal{D}$ gives rise to two profunctors:
\begin{align*}
  \mathcal{D}(-, F(-)) &: \mathcal{C} \nrightarrow \mathcal{D},
  \qquad (d, c) \mapsto \mathcal{D}(d, F(c)), \\
  \mathcal{D}(F(-), -) &: \mathcal{D} \nrightarrow \mathcal{C},
  \qquad (c, d) \mapsto \mathcal{D}(F(c), d).
\end{align*}

\subsection{Composition of Profunctors}

Two profunctors $P : \mathcal{C} \nrightarrow \mathcal{D}$ and
$Q : \mathcal{D} \nrightarrow \mathcal{E}$ can be composed. The composite
$Q \circ P : \mathcal{C} \nrightarrow \mathcal{E}$ is the \term{coend}:
\[
  (Q \circ P)(e, c) = \int^{d \in \mathcal{D}} Q(e, d) \times P(d, c).
\]
(Coends are introduced formally in Volume III; the intuition is a ``tensor
product'' over $\mathcal{D}$: matching outputs of $P$ with inputs of $Q$.)

This composition makes profunctors the morphisms in the \term{bicategory of
profunctors} $\mathbf{Prof}$. The identity profunctor on $\mathcal{C}$ is
the hom-functor $\mathcal{C}(-, -) : \mathcal{C}^{op} \times \mathcal{C} \to \mathbf{Set}$.

\subsection{Representable Profunctors}

A profunctor $P : \mathcal{C} \nrightarrow \mathcal{D}$ is \term{representable}
if it is isomorphic to $\mathcal{D}(-, F(-))$ for some functor
$F : \mathcal{C} \to \mathcal{D}$. Profunctors generalize functors by dropping
the representability requirement.

The profunctor perspective is powerful: adjunctions between $F \dashv G$ are
exactly profunctor isomorphisms $\mathcal{D}(F(-), -) \cong \mathcal{C}(-, G(-))$.

% ─────────────────────────────────────────────────────────────────────────────
\section{Exercises}

\begin{exercisesbox}
\begin{qa}
\qaitem{What does it mean for a functor to preserve a property $P$?}
       {Whenever $f$ has property $P$ in the source, $F(f)$ has property $P$ in
        the target.}

\qaitem{What does it mean for a functor to reflect a property?}
       {Whenever $F(f)$ has property $P$ in the target, $f$ has property $P$ in
        the source.}

\qaitem{Do all functors preserve isomorphisms?}
       {Yes. If $f \circ g = \operatorname{id}$ and $g \circ f = \operatorname{id}$,
        then $F(f) \circ F(g) = \operatorname{id}$ and $F(g) \circ F(f) = \operatorname{id}$.}

\qaitem{What is a faithful functor?}
       {One where each function on hom-sets $F_{a,b}$ is injective.}

\qaitem{What is a full functor?}
       {One where each $F_{a,b}$ is surjective: every morphism between images
        lifts to a morphism in the source.}

\qaitem{What is an equivalence of categories?}
       {A functor that is fully faithful and essentially surjective.}

\qaitem{Do fully faithful functors reflect isomorphisms?}
       {Yes. If $F(f)$ is an isomorphism then $f$ is an isomorphism.}

\qaitem{What is a reflective subcategory?}
       {A full subcategory whose inclusion has a left adjoint (the reflector).}

\qaitem{What is a profunctor from $\mathcal{C}$ to $\mathcal{D}$?}
       {A functor $P : \mathcal{D}^{op} \times \mathcal{C} \to \mathbf{Set}$.}

\qaitem{What notation is used for a profunctor from $\mathcal{C}$ to $\mathcal{D}$?}
       {$P : \mathcal{C} \nrightarrow \mathcal{D}$ with a slashed arrow.}

\qaitem{What is the identity profunctor on $\mathcal{C}$?}
       {The hom-functor $\mathcal{C}(-,-) : \mathcal{C}^{op} \times \mathcal{C} \to \mathbf{Set}$.}

\qaitem{How does a functor $F : \mathcal{C} \to \mathcal{D}$ give a profunctor?}
       {Two ways: $\mathcal{D}(-,F(-))$ (a profunctor $\mathcal{C} \nrightarrow \mathcal{D}$)
        or $\mathcal{D}(F(-),-)$ (a profunctor $\mathcal{D} \nrightarrow \mathcal{C}$).}

\qaitem{When is a profunctor representable?}
       {When it is isomorphic to $\mathcal{D}(-,F(-))$ for some functor $F$.}

\qaitem{How are adjunctions expressed as profunctors?}
       {$F \dashv G$ iff there is an isomorphism of profunctors
        $\mathcal{D}(F(-), -) \cong \mathcal{C}(-, G(-))$.}

\qaitem{What is the composition of profunctors $P : \mathcal{C} \nrightarrow \mathcal{D}$
        and $Q : \mathcal{D} \nrightarrow \mathcal{E}$?}
       {$(Q \circ P)(e, c) = \int^d Q(e,d) \times P(d,c)$, a coend over $\mathcal{D}$.}

\qaitem{What is the bicategory $\mathbf{Prof}$?}
       {Its objects are categories; morphisms are profunctors; 2-cells are natural
        transformations between profunctors; composition is via coend.}

\qaitem{What is RAPL?}
       {Right Adjoints Preserve Limits: if $G$ is a right adjoint then $G$ preserves
        all limits that exist in its source category.}

\qaitem{What is LAPC?}
       {Left Adjoints Preserve Colimits: if $F$ is a left adjoint then $F$ preserves
        all colimits.}

\qaitem{Give an example of a functor that does not preserve monomorphisms.}
       {The abelianization functor $\mathbf{Grp} \to \mathbf{Ab}$: injective group
        homomorphisms need not remain injective after abelianization.}

\qaitem{What does ``faithful'' mean for a forgetful functor?}
       {The underlying set function of a morphism determines the morphism; distinct
        structure-preserving maps have distinct underlying functions.}

\qaitem{Is the forgetful functor $\mathbf{Ab} \to \mathbf{Set}$ faithful?}
       {Yes: group homomorphisms with the same underlying function are equal.}

\qaitem{Is the forgetful functor $\mathbf{Ab} \to \mathbf{Set}$ full?}
       {No: not every function between abelian groups is a homomorphism.}
\end{qa}
\end{exercisesbox}

\begin{summarybox}
Functors preserve all isomorphisms but need not preserve monomorphisms, epimorphisms,
limits, or colimits unless they have adjoint structure (RAPL/LAPC). A functor is
\textbf{faithful} if injective on hom-sets; \textbf{full} if surjective; an
\textbf{equivalence} if fully faithful and essentially surjective. Fully faithful
functors reflect isomorphisms.

A \textbf{profunctor} $P : \mathcal{C} \nrightarrow \mathcal{D}$ is a functor
$\mathcal{D}^{op} \times \mathcal{C} \to \mathbf{Set}$. Every functor gives two
profunctors; adjunctions are profunctor isomorphisms. Profunctors form the
bicategory $\mathbf{Prof}$, with coend-composition.
\end{summarybox}

\section{Formal Restatement}

\begin{formalrestatement}
\textbf{Faithful/full/ff.}
$F : \mathcal{C} \to \mathcal{D}$ is faithful (resp.\ full, fully faithful) iff
$F_{a,b} : \mathcal{C}(a,b) \to \mathcal{D}(Fa,Fb)$ is injective (resp.\ surjective,
bijective) for all $a,b$.

\textbf{Reflection of isos.} $F$ fully faithful $\Rightarrow$ $F$ reflects
isomorphisms.

\textbf{Equivalence.} $F$ is an equivalence iff fully faithful and essentially
surjective; iff it has a quasi-inverse $G$ with $GF \cong \operatorname{Id}$ and
$FG \cong \operatorname{Id}$.

\textbf{Reflective subcategory.}
$\mathcal{S} \subseteq \mathcal{D}$ is reflective iff $\iota : \mathcal{S} \hookrightarrow
\mathcal{D}$ has a left adjoint $R$.

\textbf{Profunctor.}
$P : \mathcal{C} \nrightarrow \mathcal{D}$ is $P : \mathcal{D}^{op} \times \mathcal{C} \to \mathbf{Set}$.
Identity: $\operatorname{Id}_\mathcal{C} = \mathcal{C}(-,-)$.
Composition: $(Q \circ P)(e,c) = \int^d Q(e,d) \times P(d,c)$.
$F \dashv G$ iff $\mathcal{D}(F-,-) \cong \mathcal{C}(-,G-)$ as profunctors.
\end{formalrestatement}
